% =====================================================================
% Snippet: multi-zone-palette.tex
% =====================================================================
% USAGE: Reference palette for multi-zone academic figures (light bg).
% This file is a STYLE REFERENCE — copy the color definitions block
% into your figure.tex preamble. The demo shows what each color looks
% like in three contexts:
%   1. Zone background fill (light, opacity ~0.3)
%   2. Box fill (medium saturation)
%   3. Line/text/border (dark)
%
% Use rules:
%   - 1 hero/dominant color (acaPurple OR acaOrange — accent)
%   - 2-3 main color (acaBlue, acaGreen)
%   - 1-2 supporting color (acaTeal, acaGold)
%   - acaGrey for neutral/annotation
%
% 6 colors total is sweet spot. >7 = chaos. <4 = monotone.
% =====================================================================

\documentclass[tikz,border=10pt]{standalone}
\usepackage{xcolor,tikz}
\usetikzlibrary{positioning,calc}

% --- CANONICAL PALETTE (copy this block to preamble) ---
\definecolor{acaBlueLine}{HTML}{3060A8}
\definecolor{acaBlueFill}{HTML}{6890C8}
\definecolor{acaBlueBg}{HTML}{DCE6F4}

\definecolor{acaPurpleLine}{HTML}{6E3FA0}
\definecolor{acaPurpleFill}{HTML}{A878D0}
\definecolor{acaPurpleBg}{HTML}{EDE3F6}

\definecolor{acaOrangeLine}{HTML}{D06020}
\definecolor{acaOrangeFill}{HTML}{F08850}
\definecolor{acaOrangeBg}{HTML}{F9E0D0}

\definecolor{acaGreenLine}{HTML}{389050}
\definecolor{acaGreenFill}{HTML}{6CB088}
\definecolor{acaGreenBg}{HTML}{DBEDDC}

\definecolor{acaTealLine}{HTML}{1F8F8F}
\definecolor{acaTealFill}{HTML}{5AB8B8}
\definecolor{acaTealBg}{HTML}{D3EBEB}

\definecolor{acaGoldLine}{HTML}{B08820}
\definecolor{acaGoldFill}{HTML}{D8B860}
\definecolor{acaGoldBg}{HTML}{F4E8CC}

\definecolor{acaGreyLine}{HTML}{606060}
\definecolor{acaGreyFill}{HTML}{A0A0A0}
\definecolor{acaGreyBg}{HTML}{EFEFEF}

\begin{document}
\begin{tikzpicture}

% --- DEMO: 6 colors × 3 contexts ---
% Note: avoid using \fill / \line as foreach variables (they collide
% with TikZ commands). Use \colBg / \colFill / \colLine instead.
\foreach \i/\name/\colBg/\colFill/\colLine in {%
  0/Blue/acaBlueBg/acaBlueFill/acaBlueLine,%
  1/Purple/acaPurpleBg/acaPurpleFill/acaPurpleLine,%
  2/Orange/acaOrangeBg/acaOrangeFill/acaOrangeLine,%
  3/Green/acaGreenBg/acaGreenFill/acaGreenLine,%
  4/Teal/acaTealBg/acaTealFill/acaTealLine,%
  5/Gold/acaGoldBg/acaGoldFill/acaGoldLine%
} {
  \pgfmathsetmacro\yy{-\i * 1.5}
  % Zone bg (large, light)
  \fill[\colBg] (0, \yy) rectangle (3, \yy + 1.2);
  \node[anchor=north west, font=\tiny, color=\colLine]
    at (0.1, \yy + 1.1) {Zone bg (opacity 0.5)};
  % Box fill (medium)
  \node[fill=\colFill, rounded corners=2pt, font=\small\bfseries, color=white,
        minimum width=2cm, minimum height=0.8cm]
    at (4.5, \yy + 0.6) {\name Box};
  % Line/text
  \draw[\colLine, line width=1.5pt] (6.5, \yy + 0.3) -- (8, \yy + 0.9);
  \node[font=\small\bfseries, color=\colLine]
    at (9, \yy + 0.6) {\name Line};
}

\node[anchor=south, font=\bfseries] at (4.5, 1.5)
  {Canonical Light-Background Palette (use ≤6 colors)};

\end{tikzpicture}
\end{document}
